\section{Actor-Critic 数据流水线}

\subsection{采集、存储与分发}

Actor-Critic 框架的数据流可以拆成四个阶段：采集（rollout）、存储（buffer）、处理（advantage estimation）、分发（learning step）。每个阶段都有陷阱。例如，采集阶段如果数据全部来自一条 deterministic policy，会产生 low diversity，导致 replay buffer 中出现大量 near-duplicate transitions，降低实际更新的有效样本数。

\begin{knowledgebox}{Rollout 与 buffer 的协同}
为了保持采集多样性，通常采用 epsilon/entropy noise、随机策略混合等技巧；同时尽量避免把同一个轨迹重复写入 buffer，至少在 enqueue 时做 tenure check（记录 transition 被访问次数）。
\end{knowledgebox}

\subsection{优势估计与归因分配}

优势估计（advantage estimation）是 Actor-Critic 的中心。常用的算法为 GAE($\lambda$)，它在 bias-variance 之间做 trade-off：

$$\hat{A}_t = \sum_{l=0}^{T-t}(\gamma \lambda)^l \delta_{t+l},\quad \delta_t = r_t + \gamma V(s_{t+1}) - V(s_t)$$

选择 $\lambda=0.95$ 是常见的默认值，同时需要保持值函数 $\hat{V}$ 的训练稳定，否则 advantage 的方差会暴涨。

\subsection{本章小结}

Actor-Critic 的数据流水线覆盖 rollout、buffer、advantage estimation、update step。每个环节都需要监控，以防剧烈偏移带来的 instability。

